\documentclass[11pt,a4paper]{article}
\usepackage[utf8]{inputenc}
\usepackage[T1]{fontenc}
\usepackage{lmodern}
\usepackage[french]{babel}
\usepackage{amsmath,amssymb,mathtools}
\usepackage{icomma}
\usepackage{xcolor}
\usepackage{tikz}
\usepackage{pgfplots}
\usepackage[most]{tcolorbox}
\usepackage{enumitem}
\usepackage{fancyhdr}
\usepackage[hidelinks]{hyperref}
\usepackage{titlesec}
\usepackage[a4paper,margin=2.2cm,headheight=26pt,headsep=10pt]{geometry}
\usetikzlibrary{arrows.meta,positioning,calc,intersections,angles,quotes}
\usepgfplotslibrary{fillbetween}
\pgfplotsset{compat=1.18}
\definecolor{primary}{RGB}{30,75,145}
\definecolor{primarydark}{RGB}{20,50,100}
\definecolor{defcol}{RGB}{0,114,178}
\definecolor{thmcol}{RGB}{0,135,110}
\definecolor{excol}{RGB}{0,130,85}
\definecolor{remarkcol}{RGB}{120,70,180}
\definecolor{attncol}{RGB}{200,40,55}
\definecolor{methcol}{RGB}{85,85,100}
\definecolor{areacol}{RGB}{120,160,230}
\definecolor{areacol2}{RGB}{230,150,80}
\newcommand{\dd}{\mathrm{d}}
\newcommand{\R}{\mathbb{R}}
\newcommand{\ua}{\,\text{u.a.}}
\newcommand{\tg}{\tan}\newcommand{\cotg}{\cot}
\tcbset{boxbase/.style={enhanced,breakable,boxrule=0.9pt,arc=2.5pt,left=3mm,right=3mm,top=2.3mm,bottom=2.3mm,fonttitle=\bfseries,coltitle=white,toptitle=0.6mm,bottomtitle=0.6mm}}
\newtcolorbox{methode}[1][]{boxbase,colback=methcol!7,colframe=methcol,sharp corners,title={\textsc{Comment faire}\ifx\relax#1\relax\relax\else\textmd{ : \itshape #1}\fi}}
\newtcolorbox{exemple}[1][]{boxbase,colback=excol!5,colframe=excol,title={\textsc{Exemple}\ifx\relax#1\relax\relax\else\textmd{ : \itshape #1}\fi}}
\newtcolorbox{idee}[1][]{boxbase,colback=defcol!5,colframe=defcol,title={\textsc{Idée clé}\ifx\relax#1\relax\relax\else\textmd{ : \itshape #1}\fi}}
\newtcolorbox{remarque}[1][]{boxbase,colback=remarkcol!6,colframe=remarkcol,title={\textsc{Remarque}\ifx\relax#1\relax\relax\else\textmd{ : \itshape #1}\fi}}
\newtcolorbox{attention}[1][]{boxbase,colback=attncol!6,colframe=attncol,title={\textsc{Attention}\ifx\relax#1\relax\relax\else\textmd{ : \itshape #1}\fi}}
\newtcolorbox{exo}[1][]{boxbase,colback={primary!4},colframe=primary,title={\textsc{Exercice #1}}}
\newtcolorbox{corr}[1][]{boxbase,colback={thmcol!5},colframe=thmcol,title={\textsc{Corrigé #1}}}
\tikzset{axe/.style={->,>=Stealth,black!65,line width=0.7pt},courbe/.style={primary,line width=1.3pt},courbe2/.style={areacol2!90!black,line width=1.3pt},dvert/.style={gray,dashed,line width=0.7pt}}
\pagestyle{fancy}\fancyhf{}
\fancyhead[L]{\small\textcolor{primarydark}{\textbf{Fiche 7} — Intégrale d'une fonction continue}}
\fancyhead[R]{\small\textcolor{primarydark}{\textsc{Thème 4 — Intégrale définie}}}
\fancyfoot[C]{\small\thepage}
\renewcommand{\headrulewidth}{0.8pt}
\titleformat{\section}{\large\bfseries\color{primarydark}}{\thesection}{0.6em}{}

\begin{document}

\section*{Tutoriel — Thème 4 : l'intégrale définie et la formule de Newton--Leibniz}

Ce thème est le \textbf{pont} entre la primitive (une fonction) et l'intégrale définie (un
\emph{nombre}). On y apprend à passer de $F$ au résultat numérique $F(b)-F(a)$, à manier la
notation crochet, à ne jamais se tromper sur l'ordre des bornes, et à exploiter les propriétés
(Chasles, linéarité, positivité, comparaison). \emph{La mécanique du $\dd U$ et l'intégration par
parties sont détaillées dans les autres thèmes} : ici, les primitives restent simples.

\begin{idee}[la formule de Newton--Leibniz]
Soit $f$ continue sur $[a,b]$ et $F$ \textbf{une} primitive quelconque de $f$ ($F'=f$). Alors
\[
\boxed{\;\int_a^b f(x)\,\dd x=\bigl[F(x)\bigr]_a^b=F(b)-F(a).\;}
\]
Le résultat est un \textbf{nombre} (une aire algébrique, exprimée en u.a.).
\end{idee}

\subsection*{1. Pourquoi la constante $C$ disparaît-elle ?}

On peut choisir \emph{n'importe quelle} primitive, car la constante d'intégration s'élimine dans
la soustraction :
\[
\bigl[F(x)+C\bigr]_a^b=\bigl(F(b)+C\bigr)-\bigl(F(a)+C\bigr)=F(b)-F(a).
\]
Inutile donc d'écrire $+C$ dans un calcul d'intégrale définie : on prend la primitive la plus
simple (celle avec $C=0$).

\subsection*{2. La notation crochet $[\,F(x)\,]_a^b$}

Le crochet $\bigl[F(x)\bigr]_a^b$ est une \textbf{abréviation} de $F(b)-F(a)$ : on écrit d'abord la
primitive entre crochets avec les bornes, puis on substitue la borne \emph{du haut} moins la borne
\emph{du bas}. C'est le geste central de tout le thème.

\subsection*{3. L'ordre des bornes : la source d'erreur $n^{o}\,1$}

\begin{attention}[inverser les bornes change le signe]
\[
\int_a^b f(x)\,\dd x=-\int_b^a f(x)\,\dd x,\qquad\qquad \int_a^a f(x)\,\dd x=0.
\]
Une intégrale d'une borne vers \emph{elle-même} est nulle (aucune « largeur »). Et permuter le haut
et le bas ne fait que \textbf{changer le signe} du résultat.
\end{attention}

\subsection*{4. Les propriétés à connaître}

\begin{idee}[boîte à outils des intégrales définies]
Soient $f,g$ continues et $k\in\R$.
\begin{itemize}[leftmargin=1.5em,topsep=2pt,itemsep=2pt]
  \item \textbf{Chasles} : $\displaystyle\int_a^b f+\int_b^c f=\int_a^c f$ \quad(on \emph{découpe} une intégrale).
  \item \textbf{Linéarité (somme)} : $\displaystyle\int_a^b(f+g)=\int_a^b f+\int_a^b g$.
  \item \textbf{Linéarité (constante)} : $\displaystyle\int_a^b k\,f=k\int_a^b f$.
  \item \textbf{Positivité} : si $f\geqslant 0$ sur $[a,b]$ ($a<b$), alors $\displaystyle\int_a^b f\geqslant 0$.
  \item \textbf{Croissance / comparaison} : si $f\leqslant g$ sur $[a,b]$, alors $\displaystyle\int_a^b f\leqslant\int_a^b g$.
  \item \textbf{Inégalité de la moyenne} : $\displaystyle\Bigl|\int_a^b f\Bigr|\leqslant\int_a^b|f|$.
\end{itemize}
\end{idee}

Ces propriétés permettent souvent de \textbf{conclure sans tout calculer} : établir le signe d'une
intégrale, comparer deux intégrales, ou simplifier une expression avant de dériver quoi que ce soit.

\subsection*{5. Exemple type : de la primitive au nombre}

\begin{exemple}[le calcul de référence $\int_{16}^{25} x^3\,\dd x$]
\textbf{1. Primitive.} $\displaystyle\int x^3\,\dd x=\frac{x^4}{4}$ (forme puissance, $n=3$).

\textbf{2. Newton--Leibniz.}
\[
\int_{16}^{25} x^3\,\dd x=\left[\frac{x^4}{4}\right]_{16}^{25}
=\frac{25^4}{4}-\frac{16^4}{4}
=\frac{390\,625-65\,536}{4}=\frac{325\,089}{4}\approx 81\,272\ua.
\]
On substitue la borne du haut ($25$) moins celle du bas ($16$) : le résultat est un nombre.
\end{exemple}

\subsection*{6. Exemple avec compensation d'une constante et bornes}

\begin{exemple}[$\int_0^{\pi/6}\sin(2x)\,\dd x$ : faire apparaître le $\dd U$]
Une primitive de $\sin(2x)$ n'est pas $-\cos(2x)$ mais $-\tfrac12\cos(2x)$ (dérivée du composé).
On fait apparaître le facteur $2$ manquant :
\[
\int_0^{\pi/6}\sin(2x)\,\dd x=\frac12\int_0^{\pi/6}\sin(2x)\cdot 2\,\dd x
=\frac12\Bigl[-\cos(2x)\Bigr]_0^{\pi/6}
=\frac12\bigl(-\cos\tfrac{\pi}{3}+\cos 0\bigr).
\]
\[
=\frac12\Bigl(-\frac12+1\Bigr)=\frac12\cdot\frac12=\boxed{\ \frac14\ }\ua.
\]
Les angles sont en radians. On aurait aussi pu poser $U=2x$ et adapter les bornes.
\end{exemple}

\begin{remarque}[adapter les bornes lors d'un changement de variable]
Si l'on pose $U=g(x)$, on peut soit garder les bornes en $x$ et revenir à $x$ à la fin, soit
\textbf{convertir les bornes} en valeurs de $U$. Par exemple pour $\int_0^1 x\,e^{x^2+1}\,\dd x$ avec
$U=x^2+1$ : $x=0\Rightarrow U=1$ et $x=1\Rightarrow U=2$, d'où
$\tfrac12\int_1^2 e^U\,\dd U=\tfrac12(e^2-e)$. Les deux méthodes donnent le même nombre.
\end{remarque}

\end{document}
